\section{Evaluation}
\label{sec:evaluation}

We evaluate PQ-Aliro along three complementary axes: the raw communication
overhead introduced by post-quantum cryptographic primitives, the resulting
NFC transmission latency, and the security-level vs.\ overhead trade-off
across all nine parameter combinations.  All byte counts are derived from
static analysis of the implemented APDU message structures; no physical
testbed measurement is required for the core overhead figures, though we
identify empirical validation as future work.

%-----------------------------------------------------------------------
\subsection{Communication Overhead Analysis}
\label{sec:eval:overhead}

\paragraph{Methodology.}
We adopt the following counting conventions throughout this section.
(1)~\emph{Application-layer bytes} denote the TLV payload body, excluding
APDU command headers and status words.
(2)~\emph{Link bytes} add the per-frame framing cost: 6\,B per command
APDU (CLA\,$\|$\,INS\,$\|$\,P1\,$\|$\,P2\,$\|$\,Lc\,$\|$\,Le) and 2\,B
status word per response APDU (SW1\,$\|$\,SW2, appended by Android HCE
automatically).
(3)~NFC physical-layer overhead—PCB octets, CRC, and SOF/EOF framing—is
not counted, so the figures represent an \emph{upper bound on useful
payload} at the ISO-DEP layer~\cite{CSA:Aliro10}.
(4)~The LOADCERT payload size is an estimate; the actual value varies
$\pm$10\% with the X.500 Distinguished Name lengths in the Reader
certificate.
(5)~The AES-GCM IV (12\,B) is derived from a counter and is never
transmitted on the wire, so it does not contribute to link bytes.

A notable source of per-field overhead is the BER-TLV long-form length
encoding.  Any field whose encoded length exceeds 255\,B requires a
three-byte length header (\texttt{0x82}\,$\|$\,\textit{len\_hi}\,$\|$\,\textit{len\_lo})
instead of one byte, contributing 4\,B of tag-length overhead per field.
Since every PQC object—Kyber public keys ($\ge 800$\,B), ciphertexts
($\ge 768$\,B), and Dilithium public keys ($\ge 1312$\,B) and signatures
($\ge 2420$\,B)—exceeds this threshold, each such field incurs this
penalty.

\paragraph{Default configuration: ML-DSA-65 $\times$ ML-KEM-768.}
Table~\ref{tab:overhead_default} breaks down the communication overhead for
the default parameter set (Dilithium3 $\times$ Kyber768, NIST Security
Level~3~\cite{NIST:FIPS203,NIST:FIPS204}) across the six logical
message exchanges of the Expedited flow.

\begin{table}[t]
\centering
\caption{Per-phase communication overhead for the default PQ-Aliro
configuration (ML-DSA-65 $\times$ ML-KEM-768, NIST Level~3).
Chunk sizes: Reader\,$\le$\,200\,B/APDU, UD\,$\le$\,240\,B/APDU.}
\label{tab:overhead_default}
\begin{tabular}{lrrr}
\hline
\textbf{Message} & \textbf{APDUs} & \textbf{App.\ bytes} & \textbf{Link bytes} \\
\hline
AUTH0 Command   &  7   & 1{,}250   & 1{,}292   \\
AUTH0 Response  &  5   & 1{,}126   & 1{,}136   \\
LOADCERT Command & $\approx$22 & $\approx$5{,}500 & $\approx$5{,}632 \\
LOADCERT Response &  1  &       2   &       4   \\
AUTH1 Command   & 17   & 3{,}316   & 3{,}418   \\
AUTH1 Response  & 23   & 5{,}289   & 5{,}335   \\
\hline
\textbf{Total}  & $\mathbf{\approx 75}$ & $\mathbf{\approx 16{,}483}$ & $\mathbf{\approx 16{,}817 \approx 16.4\,KB}$ \\
\hline
\end{tabular}
\end{table}

\paragraph{Comparison with ECC Aliro.}
The classical ECC Aliro flow (Curve P-256) transfers approximately 1.2\,KB
per transaction: an EC ephemeral public key of 65\,B, an ECDSA signature of
roughly 72\,B, and the long-term certificate.  PQ-Aliro (D3\,$\times$\,K768)
incurs approximately $16.4/1.2 \approx \mathbf{14\times}$ more link bytes.
The dominant contributors are the Dilithium3 public key (1{,}952\,B,
$\approx$30$\times$ the 65\,B EC key) and the Dilithium3 signature
(3{,}309\,B, $\approx$46$\times$ the 72\,B ECDSA signature); the Kyber768
ciphertext (1{,}088\,B) is approximately 17$\times$ larger than an ECDH
ephemeral point.  The single largest contributing phase is LOADCERT, which
alone accounts for approximately 5.5\,KB ($\approx$33\% of the total), as
it must transmit the Reader's complete Dilithium certificate.

\paragraph{LOADCERT elimination.}
When the User Device already caches the Reader's certificate (via a
\texttt{key\_slot} index or pre-provisioning), the LOADCERT exchange can be
skipped entirely, saving approximately 22\,APDUs and 5.5\,KB.  Total
overhead drops to approximately 52\,APDUs / 11.2\,KB—a 32\% reduction
that largely closes the gap with the D2 configurations while maintaining
NIST Level~3 security.

\paragraph{All parameter combinations.}
Table~\ref{tab:overhead_configs} summarises the nine ML-DSA/ML-KEM parameter
combinations.  Within each Dilithium tier, Kyber variants account for a
modest 0.7–1.7\,KB spread; crossing Dilithium tiers (D2 $\to$ D3 $\to$ D5)
introduces 4–5\,KB jumps, driven entirely by the growth in Dilithium public
key and signature sizes~\cite{NIST:FIPS204}.

\begin{table}[t]
\centering
\caption{Communication overhead for all nine PQC parameter combinations
(with LOADCERT).  Security levels follow NIST PQC standards~\cite{NIST:FIPS203,NIST:FIPS204}.}
\label{tab:overhead_configs}
\begin{tabular}{lccc}
\hline
\textbf{Configuration} & \textbf{NIST Level} & \textbf{Total APDUs} & \textbf{Total bytes} \\
\hline
ML-DSA-44 $\times$ ML-KEM-512  & 2 &  55 & $\approx$11.9\,KB \\
ML-DSA-44 $\times$ ML-KEM-768  & 2 &  58 & $\approx$12.6\,KB \\
ML-DSA-44 $\times$ ML-KEM-1024 & 2 &  62 & $\approx$13.5\,KB \\
ML-DSA-65 $\times$ ML-KEM-512  & 3 &  72 & $\approx$15.8\,KB \\
\textbf{ML-DSA-65 $\times$ ML-KEM-768} & \textbf{3} & \textbf{75} & $\mathbf{\approx}$\textbf{16.4\,KB} \\
ML-DSA-65 $\times$ ML-KEM-1024 & 3 &  79 & $\approx$17.3\,KB \\
ML-DSA-87 $\times$ ML-KEM-512  & 5 &  95 & $\approx$21.6\,KB \\
ML-DSA-87 $\times$ ML-KEM-768  & 5 &  98 & $\approx$22.3\,KB \\
ML-DSA-87 $\times$ ML-KEM-1024 & 5 & 102 & $\approx$23.2\,KB \\
\hline
\end{tabular}
\end{table}

%-----------------------------------------------------------------------
\subsection{NFC Transmission Latency}
\label{sec:eval:latency}

\paragraph{Physical-layer model.}
ISO/IEC 14443 Type~A supports four bit rates: 106, 212, 424, and
848\,kbit/s.  At 106\,kbit/s, block-level overhead and turnaround delays
reduce the effective application throughput to approximately 8\,kB/s;
higher rates scale quasi-linearly with modest additional framing
costs~\cite{CSA:Aliro10}.  Android HCE typically negotiates 424\,kbit/s
with compatible readers, making this the expected operating point for
smartphone-based access control deployments.

\paragraph{Transmission time estimates.}
Table~\ref{tab:latency} reports estimated NFC transmission times for three
representative parameter sets at three bit rates.  The values reflect APDU
link bytes only and exclude Reader-side and UD-side cryptographic computation
times.

\begin{table}[t]
\centering
\caption{Estimated NFC transmission time (link bytes only, excluding
cryptographic computation).  Values are theoretical estimates; physical
CRC/PCB framing is not included.}
\label{tab:latency}
\begin{tabular}{lccc}
\hline
\textbf{Bit rate} & \textbf{D3$\times$K768 (16.4\,KB)} & \textbf{D2$\times$K512 (11.9\,KB)} & \textbf{D5$\times$K1024 (23.2\,KB)} \\
\hline
106\,kbit/s & $\approx$2.0\,s & $\approx$1.5\,s & $\approx$2.9\,s \\
424\,kbit/s & $\approx$0.5\,s & $\approx$0.4\,s & $\approx$0.7\,s \\
848\,kbit/s & $\approx$0.25\,s & $\approx$0.20\,s & $\approx$0.36\,s \\
\hline
\end{tabular}
\end{table}

\paragraph{Cryptographic computation overhead.}
The transmission figures above must be augmented with on-device computation
time.  On a mid-range Android device, ML-DSA-65 signature generation
requires approximately 50–100\,ms and signature verification approximately
5\,ms, consistent with results reported for the Kyber/Dilithium family on
Android hardware~\cite{NIST3PQC:Lima21}.  ML-KEM encapsulation and
decapsulation each contribute a further $<$5\,ms.  Combining transmission
and computation, the total expected latency for the default D3$\times$K768
configuration at 424\,kbit/s falls in the range of \textbf{0.6–1.5\,s}.

\paragraph{User experience impact.}
By comparison, classical ECC Aliro transfers approximately 1.2\,KB and
completes in roughly 0.2\,s end-to-end at 424\,kbit/s.  PQ-Aliro therefore
introduces a 3–7$\times$ latency increase.  At the lower end (0.6\,s), the
additional delay is likely imperceptible during a normal door-tap gesture;
at the upper end (1.5\,s), it approaches the user-experience threshold for
access-control interactions, motivating the LOADCERT-elimination optimisation
described in Section~\ref{sec:eval:overhead}.  We note that all values in
Table~\ref{tab:latency} are theoretical estimates; empirical measurement on
physical NFC hardware constitutes important future work.

%-----------------------------------------------------------------------
\subsection{Security Level vs.\ Overhead Trade-off}
\label{sec:eval:tradeoff}

The nine configurations in Table~\ref{tab:overhead_configs} span three
NIST security levels and a 1.95$\times$ overhead range (11.9\,KB to
23.2\,KB).  A Pareto analysis across the security–overhead space identifies
the following deployment profiles.

\paragraph{NIST Level~2 (D2 configurations, 11.9–13.5\,KB).}
The D2$\times$K512 configuration achieves the minimum overhead of
$\approx$11.9\,KB—a 27\% reduction relative to the default—while still
meeting the NIST Level~2 security target~\cite{NIST:FIPS203,NIST:FIPS204}.
This profile is suitable for deployments where transmission latency is
critical (e.g., high-throughput turnstiles) and where the reduced margin
against future cryptanalytic advances is acceptable within the organisation's
threat model.

\paragraph{NIST Level~3 (D3 configurations, 15.8–17.3\,KB).}
The D3$\times$K768 configuration represents the Pareto-optimal balance
between security and overhead.  It provides 128-bit post-quantum security
for both the key-encapsulation and signature primitives, and its $\approx$16.4\,KB
footprint is manageable at 424\,kbit/s.  We adopt this as the \emph{default}
for PQ-Aliro~\cite{NIST:FIPS203,NIST:FIPS204}.

\paragraph{NIST Level~5 (D5 configurations, 21.6–23.2\,KB).}
The D5$\times$K1024 configuration offers the highest available security (256-bit
post-quantum), at a cost of $+$41\% overhead relative to the default.  With
an estimated 2.9\,s transmission time at 106\,kbit/s, this profile is
impractical at the slowest NFC speed but fully viable at 424\,kbit/s
($\approx$0.7\,s) and is recommended for high-assurance installations such
as data-centre access control or government facilities.

\paragraph{Summary.}
Deployers can select the parameter set that best matches their threat
model: D2$\times$K512 for latency-sensitive or resource-constrained
environments, D3$\times$K768 as a balanced default, and D5$\times$K1024
for maximum long-term security.  The \texttt{key\_slot} optimisation
(Section~\ref{sec:eval:overhead}) orthogonally reduces overhead by
$\approx$5.5\,KB across all tiers, and should be enabled whenever the
deployment architecture permits Reader-certificate pre-provisioning on
the User Device.
